\chapter{平坦模}

除另有说明外，本章所考虑的环都假定具有单位元；所有环同态都假定把单位元映为单位元。环 A 的子环指包含 A 的单位元的子环。

如果 A 是环，M 是左 A-模，U（相应地 V）是 A（相应地 M）的加法子群，回顾一下，我们用 UV 或 U. V 表示由诸乘积 uv（其中 \(u \in U\) ， \(v \in V\) ）生成的 M 的加法子群（《代数学》第八章 §6 第 1 目）。若 a 是 A 的理想，我们写 \(a^{0} = A\) 。对任意集合 E，用 \(1_{E}\)（若不致混淆则简记为 1）表示 E 到自身的恒等映射。

回顾一下，模的公理蕴含：若 E 是环 A 上的左（相应地右）模，且 1 表示 A 的单位元，则对所有 \(x \in E\) 有 1.x = x（相应地 x.1 = x）（《代数学》第二章 §1 第 1 目）。若 E 和 F 是两个左（相应地右）A-模，回顾一下，我们用 \(\operatorname{Hom}_{\mathbf{A}}(\mathbf{E}, \mathbf{F})\)（或简记为 \(\operatorname{Hom}(\mathbf{E}, \mathbf{F})\) ）表示 E 到 F 的同态组成的加法群（出处同上，§1 第 2 目）。作为记号的滥用，我们常用 0 表示仅由其单位元组成的模。

